\Askhsh{31550_SOLUTION}{1}
\begin{ylh}{1}
    \item [Ύλη από εικόνα]
\end{ylh}
ΛΥΣΗ
\begin{alist}
\item Η $f$ είναι παραγωγίσιμη στο $(0,+\infty)$ ως άθροισμα παραγωγίσιμων με $f'(x)=e^x - \frac{1}{x}$. Η $f'$ είναι παραγωγίσιμη στο $(0,+\infty)$ ως άθροισμα παραγωγίσιμων με $f''(x)=e^x + \frac{1}{x^2} > 0$ για κάθε $x \in (0,+\infty)$, οπότε η $f$ είναι κυρτή.
\item Η $f$ είναι κυρτή οπότε η $f'$ είναι γνησίως αύξουσα στο $(0,+\infty)$. Επίσης $f'(1)=e-1>0$ και $f'\left(\frac{1}{2}\right)=e^{\frac{1}{2}}-2<0$.
Από θεώρημα Bolzano για την συνεχή $f'$ στο $\left[\frac{1}{2},1\right]$ υπάρχει $x_o \in \left(\frac{1}{2},1\right)$ με $f'(x_o)=0$ και λόγω μονοτονίας της $f'$ μοναδικό.
Επίσης, αφού $f'$ είναι γνησίως αύξουσα στο $(0,+\infty)$, έχουμε ότι για κάθε $0<x<x_o$ είναι $f'(x)<f'(x_o) \Leftrightarrow f'(x)<0$ και άρα η $f$ είναι γνησίως φθίνουσα στο $(0,x_o]$, ενώ για κάθε $x>x_o$ είναι $f'(x)>f'(x_o) \Leftrightarrow f'(x)>0$ και άρα η $f$ είναι γνησίως αύξουσα στο $[x_o,+\infty)$. Συνεπώς η $f$ παρουσιάζει ένα ακριβώς ακρότατο και μάλιστα ολικό ελάχιστο σε κάποιο $x_o \in \left(\frac{1}{2},1\right)$.
\item Είναι
\[ f'(x_o)=0 \Leftrightarrow e^{x_o} - \frac{1}{x_o} = 0 \Leftrightarrow e^{x_o} = \frac{1}{x_o} \Leftrightarrow x_o = \ln \frac{1}{x_o} \Leftrightarrow x_o = -\ln x_o \]
Το ολικό ελάχιστο είναι το $f(x_o)=e^{x_o} - \ln x_o$. Δείξαμε παραπάνω ότι $e^{x_o} = \frac{1}{x_o}$ και $x_o = -\ln x_o$ οπότε $f(x_o)=e^{x_o} - \ln x_o = \frac{1}{x_o} + x_o$.
\item Είναι $\frac{1}{x_o}+x_o > 2$ αφού
\[ \frac{1}{x_o}+x_o > 2 \Leftrightarrow \frac{1}{x_o}+x_o - 2 > 0 \Leftrightarrow \frac{x_o^2 - 2x_o + 1}{x_o} > 0 \Leftrightarrow \frac{(x_o - 1)^2}{x_o} > 0 \]
που ισχύει αφού $x_o \in \left(\frac{1}{2},1\right)$ οπότε $x_o > 0$ και $x_o \neq 1$.
Συνεπώς $f(x) \geq f(x_o) > 2$ που σημαίνει ότι η εξίσωση $f(x)=2$ είναι αδύνατη.
\end{alist}